% настройка кодировки, шрифтов и русского языка
\usepackage{fontspec}
\usepackage{polyglossia}

% рабочие ссылки в документе
\usepackage{hyperref}
%\usepackage{cite}

\usepackage[style=numeric, backend=bibtex]{biblatex}
\addbibresource{phys-bib.bib}

\usepackage{csquotes}

%%% Работа с картинками
\usepackage{graphicx} % Для вставки рисунков
\setlength\fboxsep{3pt} % Отступ рамки \fbox{} от рисунка
\setlength\fboxrule{1pt} % Толщина линий рамки \fbox{}
\usepackage{wrapfig} % Обтекание рисунков текстом

% поворот страницы
\usepackage{pdflscape}

% качественные листинги кода
\usepackage{minted}
\usepackage{listings}

% отключение копирования номеров строк из листинга, работает не во всех просмотрщиках (в Adobe Reader работает)
\usepackage{accsupp}
\newcommand\emptyaccsupp[1]{\BeginAccSupp{ActualText={}}#1\EndAccSupp{}}
\let\theHFancyVerbLine\theFancyVerbLine
\def\theFancyVerbLine{\rmfamily\tiny\emptyaccsupp{\arabic{FancyVerbLine}}}

\usepackage{multicol} 
\usepackage{tabularx}

\usepackage{extsizes}
% поля по требованию НИУ ВШЭ / МИЭМ: левое 30 мм, правое 15 мм, верх/низ 20 мм
\usepackage[left=30mm,right=15mm,
top=20mm,bottom=20mm,bindingoffset=0cm]{geometry}

% дополнительные команды для таблиц
\usepackage{booktabs}

% для заголовков
\usepackage{caption} 
\usepackage{titlesec}
\usepackage[dotinlabels]{titletoc}

% разное для математики
\usepackage{amsmath, amsfonts, amssymb, amsthm, mathtools}

% водяной знак на документе, см. main.tex
\usepackage[printwatermark]{xwatermark} 